% --- ARQUIVO: estudo-de-caso.tex ---

\section{Estudo de Caso e Configuração Experimental}
\label{sec:estudo_caso}

Para avaliar as formulações propostas sob condições realistas do Ambiente de Contratação Livre (ACL) brasileiro, desenvolveu-se um \textit{pipeline} computacional para a geração estocástica de dados fundamentais. O processo foi dividido em três frentes: simulação de preços (PLD), projeção de geração física e estruturação do portfólio financeiro.

\subsection{Modelagem Estocástica de Preços (PLD)}

O Preço de Liquidação de Diferenças (PLD) é a principal variável de incerteza do modelo. Devido à inércia estrutural do sistema elétrico — resultante da capacidade de armazenamento dos reservatórios e de ciclos macroeconômicos de demanda —, o preço de um dado mês apresenta forte dependência estatística em relação ao período imediatamente anterior. Para representar essa autocorrelação, adotou-se um modelo estatístico conhecido como PAR(1) (Autorregressivo Periódico de ordem 1), calibrado com a série histórica da CCEE de 2001 a 2025.

Os preços de energia possuem uma assimetria peculiar: eles podem sofrer picos extremos em épocas de escassez, mas não podem ser negativos. Para evitar a geração de cenários com preços irreais, a modelagem foi construída sobre o logaritmo dos dados. Para cada submercado ($s$) e mês ($m$), foram extraídos três parâmetros: a média mensal ($\mu_{s,m}$), o desvio padrão ($\sigma_{s,m}$) e o grau de correlação de Pearson com o mês passado ($\rho_{s,m}$).

A simulação dos cenários ($c$) avança um mês por vez, guiada pela seguinte relação:
\begin{equation}
    Z_{s,m}^{(c)} = \rho_{s,m} Z_{s,m-1}^{(c)} + \sqrt{1 - \rho_{s,m}^2} \cdot \varepsilon_m^{(c)}
\end{equation}

Em termos práticos, essa equação define que o novo preço padronizado ($Z$) é composto por duas componentes: a inércia herdada do mês anterior e um "choque" aleatório ($\varepsilon_m^{(c)} \sim \mathcal{N}(0,1)$). Este choque é uma variável aleatória extraída de uma distribuição normal padrão, representando a imprevisibilidade de fatores exógenos (como variações climáticas abruptas) que alteram o balanço do sistema naquele mês.

Para que o modelo respeite a realidade do Sistema Interligado Nacional (SIN), o mesmo choque estocástico ($\varepsilon_m$) foi aplicado simultaneamente aos quatro submercados. Isso garante a correlação espacial do sistema, refletindo o fato de que eventos extremos de afluência costumam impactar múltiplas regiões concomitantemente. 

Por fim, os valores simulados tiveram seu logaritmo revertido ($e^X$) e passaram por um teto e piso numérico, garantindo que nenhum cenário rompa os limites regulatórios estruturais fixados pela ANEEL para o ano vigente.

\subsection{Geração Física e Acoplamento Estocástico}

A exposição da comercializadora depende intrinsecamente do volume físico de energia de que dispõe para lastro. A geração estocástica mensal de cada usina na carteira foi simulada via método de Monte Carlo, baseando-se nas médias e desvios padrão do seu histórico operativo.

Para capturar a dinâmica do mercado de energia, implementou-se um mecanismo de \textbf{acoplamento reverso} entre a geração simulada e o PLD. Definiu-se o \textit{Fator PLD} como a razão entre o preço do cenário corrente e o preço esperado (Curva \textit{Forward}). Antes do sorteio estocástico de geração, a média esperada das usinas foi condicionalmente ajustada:
\begin{itemize}
    \item \textbf{Fontes Hidráulicas:} A média esperada sofreu ajuste inversamente proporcional ao Fator PLD, modelando o fato de que cenários de preços elevados correspondem a períodos de baixa afluência hídrica e, consequentemente, menor despacho físico das usinas.
    \item \textbf{Fontes Eólicas:} O ajuste ocorreu de forma diretamente proporcional, porém amortecida (elevado a $0.5$), capturando o fenômeno sistêmico onde períodos de elevada geração eólica costumam deprimir o PLD, especialmente na região Nordeste.
    \item \textbf{Demais Fontes (ex: Solar e Térmica):} A média de geração foi mantida estritamente igual à média histórica, desvinculada do PLD, sofrendo apenas a perturbação estocástica natural baseada em seu desvio padrão individual.
\end{itemize}

\subsection{Estruturação de Portfólio e Curva Forward}

A posição inicial do agente (Carteira Legada) foi parametrizada como uma posição de Venda \textit{Short} (\textit{hedge}), cujo volume corresponde a um percentual fixo da geração física média esperada de seus ativos. O preço médio desses contratos foi balizado pelo histórico público de liquidação de contratos bilaterais MCSD de energia nova.

Para a estruturação das opções de decisão do otimizador (a "Mesa de Operações"), calcularam-se os valores esperados do PLD em cada nó, extraindo-se a \textbf{Curva Forward} do sistema. Foram então instanciados contratos de compra e venda com durações variadas ($d_i$). Para emular a liquidez e o prêmio de risco do mercado real, os preços dos ativos negociáveis foram estabelecidos aplicando-se um \textit{Spread} assimétrico sobre a Curva Forward:
\begin{align}
    K_i^{\text{Compra (Ask)}} &= \mathbb{E}[PLD_{s, d_i}] \times (1 + Spread) \\
    K_i^{\text{Venda (Bid)}} &= \mathbb{E}[PLD_{s, d_i}] \times (1 - Spread)
\end{align}
Dessa forma, o otimizador só assume novas posições físicas se o ganho na gestão de risco da árvore de cenários justificar o custo de transação (\textit{spread}) embutido nos contratos daquela maturidade.